\section{Conclusions}

The experimental results, both in simulation and on the constructed platform, demonstrate that the proposed neural architecture enables the quadruped robot to perform stable and adaptive navigation in controlled environments with different stimulus configurations and topographical variations. In scenarios featuring stimuli of different natures, the physical robot reaches an appetitive stimulus located 1.70 m ahead while avoiding obstacles and evading the aversive stimulus in an average time of less than 80~s, maintaining efficient trajectories and a constant frontal orientation toward the target. Overall, these results confirm that combining bioinspired control with multimodal sensing constitutes an effective strategy for autonomous locomotion in robotic inspection contexts. 

The observed performance highlights the effectiveness of the obstacle-perception module and the locomotion scheme controlled by central neural networks. The LiDAR-driven neuronal ring enabled early and directional detection of objects, generating inhibitory or excitatory signals depending on stimulus nature. This sensory processing, coupled to the locomotion module, facilitated trajectories with minimal deviations from the optimal path, maintaining a consistent safety distance from non-relevant obstacles. Sustained activity in neuron $L_4,$ along with the activation pattern of neuron $X_{13}$, correlated directly with persistent orientation toward the target, preventing oscillations or unnecessary course corrections. This coherence between perception and motor control confirms that integrating a directional sensory system with a central pattern generator modulated by stimulus type is sufficient to ensure stable and efficient displacement even in the absence of environmental conflict. 

In scenarios presenting simultaneous appetitive and aversive stimuli, the robot consistently prioritized the most relevant objective through the basal-ganglia-inspired arbitration module. The interaction among the STR, STN, GPe, and GPi implemented a winner-take-all mechanism, actively suppressing responses toward non-priority stimuli. During the tests, predominant activation in the channel associated with the selected stimulus effectively inhibited competing channels, reducing the probability of erratic trajectory changes. This behavior yielded direction-switching latencies below 600~ms, ensuring fast and stable responses under sensory conflict. The alignment between decision-neuron activity peaks and executed maneuvers supports the hypothesis that a basal-ganglia-inspired architecture combined with multimodal sensory processing offers clear advantages in environments where goal prioritization is critical for mission success. 

On irregular surfaces and steep slopes, the system exhibited an adaptive response grounded in the direct integration of inertial data into the locomotion control module. Continuous monitoring of roll, pitch, and the standard deviation of vertical acceleration ($\sigma_{az}$) enabled early detection of instability conditions. When the predefined thresholds for each metric were exceeded, the neural architecture immediately triggered a transition from H-mode to quadrupedal mode, characterized by a gradual activation pattern in the modulator neurons ($X_{0}$–$X_{2}$) and a controlled inhibition of the previous gait. This transition not only maintained postural stability but also significantly reduced lateral and longitudinal oscillation amplitudes, minimizing the risk of loss of balance. The relationship between IMU signal intensity and postural-response activation confirms that combining kinematic and dynamic criteria is an effective strategy for locomotor adaptation in adverse terrains, ensuring both structural integrity of the robot and continuity of the mission. Likewise, it is noted that the first version of the gait-decision network implemented in simulation produces variant transient states during transitions to and from the quadrupedal mode; therefore, exploring additional IMU-derived variables or signal-processing approaches could yield more robust terrain-state inference. 

These results are consistent with those reported by \citet{sarvestani2013computational}, who demonstrated that a basal-ganglia-based arbitration scheme can manage sensory conflicts and robustly prioritize motor behaviors. However, this work introduces several substantial contributions beyond the state of the art. First, a multimodal robot capable of switching between distinct locomotion patterns depending on terrain conditions was developed, a feature that increases operational versatility compared to unimodal platforms. Second, an inertial sensor was incorporated for real-time monitoring of postural stability, enabling early detection of disturbances and immediate transitions to more stable gaits. Finally, an original locomotor neural network ---modulated by the Naka–Rushton function--- was designed, yielding smoother transitions and reducing undesired oscillations. Together with the basal-ganglia-inspired arbitration architecture, these innovations produced more stable trajectories, demonstrating a solid computational foundation for autonomous gait arbitration in hybrid biomorphic platforms. 

Several limitations must be acknowledged. The neural and locomotion-control parameters were manually tuned and kept constant during all experiments, without applying adaptive optimization processes that could enhance performance in more diverse scenarios. Additionally, system resilience was not assessed under partial hardware failures or temporary sensory dropout-conditions likely to occur in real industrial environments. Classification metrics for strengthening the MLP-based terrain classifier were also not evaluated. Furthermore, the multimodal sensory integration does not follow a formal fusion framework in the classical signal-processing sense---such as Kalman filtering, Bayesian estimation, or deep sensor fusion. Instead, cross-modal integration emerges implicitly within the bio-inspired neural network: each sensory channel feeds differentiated neural pathways that converge through Winner-Take-All competition dynamics in the basal ganglia module, effectively distributing the fusion function across the network topology. While this approach is biologically grounded, it has not been formally compared against standard fusion methods. Similarly, visual stimulus detection relies on color-based segmentation calibrated under controlled illumination, without evaluation under varying lighting conditions or occlusion. Notwithstanding these design simplifications, the system demonstrated empirical robustness to stochastic sensory perturbations: under artificially introduced noise levels ($\sigma \in [0, 4]$), mission completion times remained tightly bounded, decision latency was stable at approximately $47$~ms, and kinematic variance was minimal across all tested conditions---evidencing that the neural architecture effectively absorbs sensory disturbances without compromising behavioral integrity. These constraints do not invalidate the results but suggest that the next development stage should focus on validating the proposed architecture in a prototype that includes robustness and adaptability tests under real, non-controlled conditions.

The findings provide a highly efficient baseline that demonstrates how lightweight, bio-inspired neural networks can successfully manage multimodal gait transitions on embedded hardware, offering significant insights for the future development of industrial-grade autonomous inspection units. The ability to switch locomotion modes based on stability metrics derived from inertial sensing expands the robot's operational range, ensuring reliable displacement on flat, rough, and inclined surfaces. The integration of a multimodal neural architecture, a basal-ganglia-inspired arbitration module, and an original locomotor network enables dynamic goal prioritization, collision avoidance, and reduced response times to unexpected environmental changes. When scaled onto robust physical platforms, this low-cost computational paradigm demonstrates how decentralized, bio-inspired architectures could optimize energy efficiency and behavioral reliability. Consequently, it serves as a powerful stepping stone toward bridging the gap between open-loop exploratory prototypes and the rigorous stabilization required by unconstrained, real-world deployment in power-plant monitoring or automated industrial operations. 

The modular nature of the architecture facilitates its adaptation to different robotic platforms and sensor types, broadening its applicability across various mission scenarios. Furthermore, for industrial inspection contexts, it is relevant to detail complementary sensing technologies ---beyond camera, LiDAR, and IMU--- to avoid reliance on superficial evidence. The current platform already includes a pipeline of inertial and vision data (including stagnation detection) feeding the decision and gait-transition modules, enabling incorporation of new sensor modalities without modifying the system’s core arbitration and sensorimotor fusion principles. 

Overall, this work demonstrates that a sensorimotor architecture inspired by basal-ganglia-like circuits ---capable of arbitrating between locomotion modes and reactive behaviors driven by visual and inertial cues--- allows a real robotic platform to navigate, stabilize, and make decisions robustly in environments with obstacles and adverse terrains. The integration of camera, LiDAR, and IMU into low-latency decision loops resulted in more stable trajectories, timely gait transitions, and accurate goal approach without collisions, validated in both simulation and physical experiments. 

Beyond mobility, the proposed framework is extensible: incorporating NDT modalities (infrared, ultrasound, hyperspectral) and fusing them over 2.5D maps will enable more reliable damage diagnostics for industrial inspection. Long-term field validation, adaptive threshold learning, and policy optimization are recommended to reduce manual tuning, along with developing safety and traceability protocols aligned with current standards. With these advances, the platform can evolve from a functional prototype into a practical autonomous inspection tool.

In conclusion, this paper presented a biologically inspired neural controller modeled after basal ganglia mechanisms for adaptive gait and mode arbitration on a hybrid wheel-legged quadruped. Although the physical prototype displays hardware-level boundaries ---including low-torque actuation and open-loop locomotion profiles--- the framework demonstrates that complex, multi-stimulus mobility decisions can be executed without heavy computational overhead or expensive training cycles. Future work will focus on migrating the arbitration logic to industrial-grade platforms, implementing closed-loop dynamic gait replanning, integrating formal sensor fusion, and validating the method in unconstrained outdoor environments.

\section*{Acknowledgements}

The authors would like to thank the Neurocontrol Research Group at the Universidad Autónoma de Occidente for their continuous technical and academic support during the development of this project.

\section*{Author Contributions}

\textbf{Samuel F. Carlos Rodriguez:} Methodology, Software, Investigation, Writing--original draft, Supervision.
\textbf{Diego A. Caro Vaca:} Software, Investigation, Writing--original draft, Supervision.
\textbf{Sebastian Grajales Pulgarín:} Methodology, Software, Investigation, Resources.
\textbf{David Caicedo Samboni:} Investigation, Formal analysis, Writing--original draft, Visualization.
\textbf{Julián Hurtado-López:} Supervision, Project administration.
\textbf{David F. Ramirez-Moreno:} Project administration, Resources.
\textbf{All authors:} Conceptualization, Writing--review \& editing, Validation.

\section*{Declaration of Competing Interest}

The authors declare that they have no known competing financial interests or personal relationships that could have appeared to influence the work reported in this paper.

\section*{AI-assisted writing disclosure}

During the writing of this manuscript and its subsequent revisions, we used artificial intelligence-based tools, including ChatGPT (OpenAI), as editorial support to improve the grammar, clarity, tone, coherence, and overall writing quality of the text. These tools were used solely for linguistic and stylistic refinement and were not used to generate scientific results, data, analyses, or conclusions. The authors critically reviewed, edited, and approved the final content and take full responsibility for the accuracy and integrity of the manuscript.

%% For citations use: 
%%       \citet{<label>} ==> Lamport [21]
%%       \citep{<label>} ==> [21]
%%
%%Example citation, See \citet{lamport94}.

%% If you have bib database file and want bibtex to generate the
%% bibitems, please use
%%
